% Mục 7 (bản rút gọn). Rút từ sections/threats.tex; nguồn từng số ghi ở đó (results/<máy>/*, machine.txt,
% scripts/report_numbers/corr_stats.py, corr_stats2.py, extra_analyze.py, rlock_analyze.py).
\emph{Construct validity} (số đo có đo đúng đại lượng muốn đo không). Cận dưới của phần lock dùng $t$ đo ở 1 thread, nên chứa chi phí lấy và nhả
lock của lock nhanh nhất. Ở $N=11$, chi phí này chênh 41\,\% giữa các lock. Cận dưới bỏ qua giới hạn 10
core khi $p > 10$ và $\mathrm{OW} = 15$ (ở 20 thread sàn thật cao hơn 1,6 lần). Cận dưới còn giả định
mọi core chạy \texttt{is\_prime} nhanh như ở lần chạy 1 thread. Khi số thread vượt số P-core, một phần
việc chạy trên E-core chậm hơn, nên một phần độ chậm không do lock (ô đối chứng ở RQ2). Phép quy về cùng
số vòng lặp giả định thời gian tỉ lệ với số vòng. Giả định này yếu nhất ở MCS quá tải (18--286 vòng mỗi
thread trên \PCname{apple-m5}, 10--56 trên \PCname{LAP1184}), ở counter \PCname{lock_mcs} quá tải
(10--200 lần tăng mỗi thread) và ở các ô barrier sụp (50--69 episode mỗi lần chạy). Ở $\mathrm{OW}=8$,
thời gian của mọi counter, kể cả \PCname{faa}, tăng ít nhất gấp đôi từ 4 lên 10 thread, và nhóm không có
phép đo đối chứng để tách phần do E-core. Phép đo khoảng trống trao lock thêm hai lần đọc đồng hồ vào mỗi
critical section, nên báo cáo chỉ dùng thống kê khoảng trống của nó. Độ công bằng so thời điểm kết thúc
của các thread, nên không thấy các đợt độc chiếm luân phiên ngắn.
% nguồn: locksingle N=11: 6,50/4,60 = 1,41; plot.aggregate mcs reduced iters; counter.csv lock_mcs iters <= 200;
% barrier.csv sr, sr_bo_spin 12-20 thread: iters 69/66/50; counter.csv OW=8 faa 10t/4t 2,13.

\emph{Internal validity} (khác biệt đo được do yếu tố được thay đổi hay do nhiễu). Không lần chạy nào ghim thread vào core, và macOS không cho ghim thread. Ở
$N=1009$, thời gian một lời gọi \texttt{is\_prime} trong 420 lần chạy 1 thread trên \PCname{apple-m5} đi
từ 468 tới 1\,713~ns. Nhóm cho rằng độ phân tán này phần lớn do thread rơi vào E-core, nhưng không tách
được ảnh hưởng của E-core với tải nền. Load average 1 phút trong đợt đo lock có trung vị khoảng 25 ở mọi
số thread, kể cả lần chạy 1 thread (12,4--51,2). Con số này cao gấp đôi số thread trung bình theo thời
gian của đợt đo lock (12,0).
% nguồn: lock.csv cột load1: 12,42..51,18, trung vị 25,4; số thread trung bình có trọng số seconds 11,97.
Nhóm đo lại năm lock ở $N = 1009$, $\mathrm{OW} = 0$, 1--10 thread, mỗi lần chạy chỉ bắt đầu khi load
average dưới 3 hoặc sau 2 phút chờ. Các đường cong giữ hình dạng khi tải thấp, còn điểm gãy của MCS ở 10
thread nhỏ đi một nửa (0,53 lần, tức chậm khoảng 7 lần cận dưới thay vì 13). Ngoài MCS ở 10 thread,
trung vị mới bằng 0,83--1,05 lần trung vị của đợt đo chính, tức lệch tới 17\,\%. Báo cáo không diễn giải
những khác biệt nhỏ hơn mức này. Lần đo lại khác đợt đo chính cả về phiên chạy, nên không tách được riêng
ảnh hưởng của tải. Vùng vượt số core chưa được đo lại.
% nguồn: extra_analyze.py (results/apple-m5/extra/lowload.csv): 0,83..1,05; mcs p=10: 0,53.
Nhóm không đọc nhiệt độ hay tần số CPU. Thời gian 1 thread không trôi quá 3,4\,\% qua 7 lượt lặp.
Xác suất 98,4\,\% tính cho thanh min--max ở mục Phương pháp đo giả định các lần lặp độc lập và cùng
phân bố, trong khi tải nền thay đổi trong đợt đo, nên thanh này chỉ là khoảng min--max của các lần lặp,
không phải khoảng tin cậy. Trung vị không thể hiện phân bố hai đỉnh, và chỉ thanh min--max cho thấy phần
này.
% nguồn: corr_stats2.py, corr_stats.py (tỉ số theo lượt 0,977..1,034).
Trên \PCname{LAP1184}, load average chỉ đếm tiến trình trong máy ảo. Tải của Windows, việc Windows xếp
lịch cho các CPU ảo và chế độ điện trong lúc đo không được ghi lại. \texttt{lscpu} báo tần số cố định,
nên không thấy được throttling. Thời gian một lời gọi ở 1 thread, $N=1009$, đi từ 1\,463 tới 3\,176~ns
qua 455 lần chạy. Nhóm không tách được phần của bộ hẹn giờ ảo với phần của timer slack trong thời gian
ngủ của \texttt{sleep\_for}. Pause-loop exiting, cơ chế hypervisor tạm dừng một CPU ảo quay vòng chờ quá
lâu, và việc đánh thức thread qua ngắt ảo là những giải thích khác, chưa đo, cho việc barrier spin chậm
ngay ở 12 thread.
% nguồn: machine.txt (cpu MHz 2496.0, virtualization: Microsoft); lock.csv LAP1184 threads=1, n=1009.

\emph{External validity} (kết quả có khái quát được ra ngoài ba máy đã đo không). Số đo đến từ hai máy Apple Silicon (Apple clang 21, macOS) và một máy x86-64
(g++ 11.5, Linux trong máy ảo WSL2), cả ba đều \texttt{-O2}. Hai máy Apple khác nhau ở số P-core và số
core, nhưng cũng ở thế hệ chip (M3 Pro, M5), dung lượng L2 mỗi cụm và phiên bản macOS. Vì vậy việc điểm
gãy dời theo số P-core phù hợp với giải thích hai loại core, nhưng không chứng minh nó. Máy x86-64 khác
hai máy Apple cùng lúc ở kiến trúc, memory model, compiler, hệ điều hành và việc có máy ảo hay không, nên
RQ6 chỉ gán một khác biệt của máy này cho một nguyên nhân khi có bằng chứng riêng. Nhóm không có máy x86
chạy Linux trực tiếp hay máy nhiều socket. \texttt{is\_prime} không chạm dữ liệu dùng chung. Critical
section thật thường ghi dữ liệu dùng chung, nên chi phí chuyển giao lock thực tế có thể lớn hơn.
